\documentclass{article}

\usepackage{geometry}
\geometry{a4paper}

\usepackage[superscript,biblabel]{cite}
\bibliographystyle{naturemag}
\usepackage[hidelinks]{hyperref}
\usepackage{authblk}
\usepackage{tabularray}
\usepackage{graphicx}
\usepackage{amssymb}

% Make affiliations run inline
\renewcommand\Affilfont{\small}
\renewcommand\Authfont{\normalsize}

\makeatletter
\renewcommand\AB@affilsepx{;\hspace{0.6em}\protect\Affilfont}
\makeatother

\renewcommand{\thetable}{S\arabic{table}}
\renewcommand{\thefigure}{S\arabic{figure}}
\renewcommand{\thesubsection}{\arabic{subsection}}

\begin{document}

\title{Population-scale Ancestral Recombination Graphs with tskit 1.0\\
\large Supplementary Information}

\input{authors.tex}

\date{}
\maketitle

\subsection*{Published tools using tskit}
We performed a search for published software using tskit.
To be included in this list, the software must be described in a publication
as a named tool intended for reuse and must use tskit.
This list is intentionally conservative, and does not include
(for example) publications depending on bespoke analysis
pipelines using tskit.
For each tool in Table S1, we list: its name (with a clickable
hyperlink to its source code repository);
the main implementation
language(s) of the tool based on GitHub summaries;
the DOI (with clickable hyperlink) of the publication in which this
tool was first described (or, in the case of SLiM, the first
publication in which tskit-based features were described);
and the year of publication and number of citations to this
publication according to data from OpenAlex
(\url{https://openalex.org}) retrieved on 2026-01-23.

\input{tools_table.tex}

\subsection*{Tskit functionality and benchmarks}
In this section we provide an overview of tskit's functionality,
how this overlaps with adjacent libraries, and summarise
performance benchmarks performed in previous publications.

To give a sense of the breadth of functionality provided by tskit,
Table S2 lists properties and operations provided by tskit
broken into broad categories.
There is no direct competitor to tskit, so a head-to-head feature
comparison is not possible; instead, we contrast tskit with three
Python libraries that target adjacent niches:
ARGneedle-lib\cite{zhang2023biobank}, which infers and analyses
ancestral recombination graphs from biobank-scale data;
the BTE interface (\url{https://github.com/jmcbroome/BTE})
for matUtils\cite{mcbroome2021daily},
a toolkit for manipulating and analysing mutation-annotated trees of SARS-CoV-2;
and DendroPy\cite{sukumaran2010dendropy}, a long-established
general-purpose phylogenetic scripting library.
For each operation we mark whether the comparison library offers a
comparable feature.
The list is illustrative rather than exhaustive and reflects
library state at the time of writing.

Tskit is the result of more than a decade of sustained development,
and the algorithmic advances it represents have been described and benchmarked in
a series of publications. The succinct tree sequence data format
and the core algorithms for sequential tree construction
were first introduced as part of the msprime coalescent
simulator in \cite{kelleher2016efficient}. [Mention the
theoretical algorithmic results.] This paper demonstrated
the efficiency of the approach in a benchmark that illustrated
that using the succinct tree sequence format and efficient
tree iteration was more than a million times faster than accessing
the trees individually using Newick and the currently available
libraries in a simulation with 100,000 individuals.

The data model was generalised to enable forwards-time simulations
in \cite{kelleher2018efficient}, and the columnar tables API.
The tskit library itself was first introduced here. The critical
``simplify'' method was described in detail and the time and space
complexity analysed rigorously. This function is the core method
required for forwards-time ARG simulations, and
is now used in varied applications~\cite{haller2018tree}.

The genealogical nearest neighbourhs (GNN) statisic was introduced
in~\cite{kelleher2019inferring}.

The general API for computing population genetics statistics
was introduced in~\cite{ralph2020efficiently}, which also introduced
the concept of ``branch based statistics''. Along with a rigorous
theoretical analysis, benchmarks against existing libraries showed
that tskit's implementation was up to 100,000 times faster.

[Summarise also: \cite{fritze2026forest}, \cite{lehmann2026on}
and \cite{lee2026genetic}. Although Lee et al 2026 doesn't describe
updates to tskit itself (I think) it has benchmarks of the GRM
against argneedle-lib.]

\begin{tblr}[
        long,
        caption={Overview of tskit's functionality, grouped into eight
            broad categories. Every operation listed is implemented in
            tskit; the columns indicate whether each comparison library
            offers a comparable feature: \checkmark{} denotes full
            support, $\circ$ denotes partial or restricted support, and
            a blank cell denotes none.
            }
     ]{
        colspec=Xccc,
        rowhead=1,
        width=\linewidth,
        colsep=3pt,
         }
\hline
Operation & ARGneedle-lib & matUtils (BTE) & DendroPy \\
\hline
\SetCell[c=4]{c} \textbf{Data model}\\
\hline[dashed]
Lossless ARG representation & \checkmark & & \\
Mutations integrated with topology & \checkmark & \checkmark & \\
Integrated reference sequence & & & \\
Individual / ploidy abstraction & $\circ$ & & \\
Population structure encoding & & & \\
Columnar (NumPy) data access & & & \\
\hline[dashed]
\SetCell[c=4]{c} \textbf{File formats}\\
\hline[dashed]
Efficient binary format & \checkmark & \checkmark & \\
VCF export & & \checkmark & \\
Newick export & \checkmark & \checkmark & \checkmark \\
FASTA export & & & \checkmark \\
NEXUS export & & & \checkmark \\
\hline[dashed]
\SetCell[c=4]{c} \textbf{Tree operations}\\
\hline[dashed]
Iterate trees along a genome & & & \\
Tree traversal (pre-/post-order) & & \checkmark & \checkmark \\
MRCA and common-ancestor queries & \checkmark & \checkmark & \checkmark \\
Branch length and total branch length & $\circ$ & \checkmark & \checkmark \\
Tree topology comparison (KC, RF) & \checkmark & & \checkmark \\
Tree balance and shape statistics & & $\circ$ & \checkmark \\
\hline[dashed]
\SetCell[c=4]{c} \textbf{ARG/tree editing}\\
\hline[dashed]
Simplify (sample-restricted history) & & & \\
Subset by sample or clade & & \checkmark & \checkmark \\
Union of tree sequences & & & \\
Keep / delete genomic intervals & $\circ$ & & \\
Trim flanking regions & \checkmark & & \\
\hline[dashed]
\SetCell[c=4]{c} \textbf{Population-genetic statistics}\\
\hline[dashed]
Nucleotide diversity ($\pi$), segregating sites & & $\circ$ & \checkmark \\
Tajima's $D$ & & & \checkmark \\
$F_{ST}$ and divergence & & & \\
$f$-statistics ($f_2$, $f_3$, $f_4$, Patterson's $D$) & & & \\
Allele frequency spectrum (one-way and joint) & & & $\circ$ \\
Linkage disequilibrium & & & \\
Branch-mode statistics on trees & & & \\
\hline[dashed]
\SetCell[c=4]{c} \textbf{Ancestry and relatedness}\\
\hline[dashed]
IBD segment extraction & & & \\
Genealogical nearest neighbours (GNN) & & & \\
Genetic relatedness matrix & \checkmark & & \\
Pairwise divergence / coalescence times & \checkmark & $\circ$ & $\circ$ \\
Mean descendants & & $\circ$ & \\
\hline[dashed]
\SetCell[c=4]{c} \textbf{Mutations and variants}\\
\hline[dashed]
Variant / genotype iteration & $\circ$ & \checkmark & \\
Haplotype reconstruction & & \checkmark & \\
Mutation placement / parsimony & \checkmark & \checkmark & \\
\hline[dashed]
\SetCell[c=4]{c} \textbf{Visualisation}\\
\hline[dashed]
SVG tree drawing & & & \\
Tree-sequence (multi-tree) visualisation & & & \\
ASCII / text tree rendering & & & \checkmark \\
\hline[dashed]
\SetCell[c=4]{c} \textbf{Metadata and provenance}\\
\hline[dashed]
Structured metadata with schemas & & $\circ$ & \\
Provenance recording and validation & & & \\
\hline
\end{tblr}


\bibliography{paper}

\end{document}
